\Qhead{5.1}{Антенны}{Вибратор Герца. Ближняя и дальняя зоны. Диаграмма направленности. Коэффициент усиления и коэффициент рассеяния антенны. Антенны для ДВ-, СВ- и СВЧ-диапазонов. Параболическая антенна. Фазированные антенные решетки. Эффективная площадь и шумовая температура приемной антенны.}
\Subfile{5.1.1}{sections/05_antennas/q01/s01}
\Subfile{5.1.2}{sections/05_antennas/q01/s02}
\Subfile{5.1.3}{sections/05_antennas/q01/s03}
\Subfile{5.1.4}{sections/05_antennas/q01/s04}
\Subfile{5.1.5}{sections/05_antennas/q01/s05}

\recap{Антенна преобразует ток в пространственное поле и обратно. Направленность описывает геометрию излучения, усиление учитывает потери, эффективная площадь — приём. Поляризация, согласование и шум окружающей среды определяют реальную чувствительность.}
\board{Дипольная диаграмма; $D,G,A_e$; множитель решётки; $P_n=k_BT_A\Delta f$.}
\oral{Почему рост направленности не обязательно уменьшает шумовую температуру антенны?}
\bait{Можно оставить тёплые боковые лепестки: они собирают шум Земли даже при наведении главного луча на холодное небо.}
\sources{Коэффициент рассеяния: Шарков~\cite{sharkov}, §5.4, с.~184; Сорокин~\cite{sorokinantenna}, PDF-с.~4--15; Биленко и соавт.~\cite{bilenko}, §§2.3.3--2.3.4. Формулы Герца, решётки и шумового взвешивания получены в явно указанных нормировках; неоднозначность «коэффициента рассеяния» сохранена.}
